\subsection{Discrete approximation for compact domains}\label{subsec:approximation-discrete}
\paraid{p-39f5655ab8f2}
We now construct the approximating maps and prove discreteness at
every point of the ambient Gromov--Hausdorff space.

\paraid{p-db4cf11c2a9d}
We first choose a positive Lipschitz scale adapted to an open cover.

\begin{lemma}\label{lem:cover-adapted-scale}
\paraid{p-90d3aa515184}
Let
$\OpenCover$
be an open cover of
$\GHSpace$.
Then there exists a
$1/16$-Lipschitz function
$\CoverScale\colon\GHSpace\to(0,1/16]$
with the following property.
For every
$\MetricSpace{\BaseCarrier}{\MetricSymbol},
 \MetricSpace{\ComparisonCarrier}{\ComparisonMetric}\in\GHSpace$,
if
\[
 \GHDistance(\MetricSpace{\BaseCarrier}{\MetricSymbol},
             \MetricSpace{\ComparisonCarrier}{\ComparisonMetric})
 <\CoverScale\MetricSpace{\BaseCarrier}{\MetricSymbol},
\]
then there exists
$U\in\OpenCover$
with
\[
 \{\MetricSpace{\BaseCarrier}{\MetricSymbol},
   \MetricSpace{\ComparisonCarrier}{\ComparisonMetric}\}\subset U.
\]
\end{lemma}

\begin{proof}
\paraid{p-8da60d9b543c}
Use the convention that distance to the empty set is
$\infty$.
For every
$\MetricSpace{\BaseCarrier}{\MetricSymbol}\in\GHSpace$,
put
\begin{equation}\label{eq:cover-scale-2}
 \CoverScale\MetricSpace{\BaseCarrier}{\MetricSymbol}
 =\frac1{16}\sup_{\OpenSubset\in\OpenCover}
   \min\left\{1,\dist_{\GHDistance}
       \bigl(\MetricSpace{\BaseCarrier}{\MetricSymbol},\GHSpace\setminus\OpenSubset\bigr)\right\}.
\end{equation}
Each truncated distance function is
$1$-Lipschitz,
and their supremum is positive because
$\OpenCover$
is an open cover.
Thus,
for every
$\MetricSpace{\BaseCarrier}{\MetricSymbol},
 \MetricSpace{\ComparisonCarrier}{\ComparisonMetric}\in\GHSpace$,
we have
\begin{equation}\label{eq:scale-lipschitz-1}
 0<\CoverScale\MetricSpace{\BaseCarrier}{\MetricSymbol}\leq\frac1{16},
\end{equation}
and
\begin{equation}\label{eq:scale-lipschitz-2}
 \abs{\CoverScale\MetricSpace{\BaseCarrier}{\MetricSymbol}
      -\CoverScale\MetricSpace{\ComparisonCarrier}{\ComparisonMetric}}
 \leq\frac1{16}\GHDistance(\MetricSpace{\BaseCarrier}{\MetricSymbol},
                          \MetricSpace{\ComparisonCarrier}{\ComparisonMetric}).
\end{equation}
Now let
$\MetricSpace{\BaseCarrier}{\MetricSymbol},
 \MetricSpace{\ComparisonCarrier}{\ComparisonMetric}\in\GHSpace$
satisfy
$\GHDistance(\MetricSpace{\BaseCarrier}{\MetricSymbol},
             \MetricSpace{\ComparisonCarrier}{\ComparisonMetric})
 <\CoverScale\MetricSpace{\BaseCarrier}{\MetricSymbol}$.
By \eqref{eq:cover-scale-2},
there exists
$\mathscr O\in\OpenCover$
such that
\[
 \GHDistance(\MetricSpace{\BaseCarrier}{\MetricSymbol},
             \MetricSpace{\ComparisonCarrier}{\ComparisonMetric})
 <\CoverScale\MetricSpace{\BaseCarrier}{\MetricSymbol}
 <\dist_{\GHDistance}
   \bigl(\MetricSpace{\BaseCarrier}{\MetricSymbol},\GHSpace\setminus\mathscr O\bigr).
\]
Both spaces therefore belong to
$\mathscr O$.
\end{proof}

\begin{theorem}\label{thm:discrete-approximation}
\paraid{p-07c6df8fb234}
Let
$\{\ApproximationDomain _\FamilyIndex \}_{\FamilyIndex \in\NonnegativeIntegers}$
be a sequence of compact metrizable spaces.
For each
$\FamilyIndex\in\NonnegativeIntegers$,
let
$\InputMap _\FamilyIndex \colon \ApproximationDomain _\FamilyIndex \to\GHSpace$
be continuous.
Let
$\OpenCover$
be an open cover of
$\GHSpace$.
Then there exist continuous maps
$\ApproximatingMap _\FamilyIndex \colon \ApproximationDomain _\FamilyIndex \to\GHSpace$
such that each
$\ApproximatingMap _\FamilyIndex $
is
$\OpenCover$-close
to
$\InputMap _\FamilyIndex $
and the family
$\{\ApproximatingMap _\FamilyIndex (\ApproximationDomain _\FamilyIndex )\}_{\FamilyIndex \in\NonnegativeIntegers}$
is discrete in
$\GHSpace$.
\end{theorem}

\begin{proof}
\paraid{p-bce29ad7f159}
We first dispose of the empty domains.
For each empty domain
$\ApproximationDomain_\FamilyIndex$,
choose
$\ApproximatingMap_\FamilyIndex$
to be the unique empty map.
Empty images do not affect closeness or discreteness,
so the assertion holds if every domain
$\ApproximationDomain_\FamilyIndex$
is empty.
We now consider only the nonempty domains and put
$\DomainIndexSet
 =\{\FamilyIndex\in\NonnegativeIntegers\mid
    \ApproximationDomain_\FamilyIndex\ne\emptyset\}$.
For each
$\FamilyIndex\in\DomainIndexSet$,
we define
a positive integer
$\EquilateralSize_{\FamilyIndex}$
recursively,
and
for
$\DomainPoint\in\ApproximationDomain_{\FamilyIndex}$,
we construct the approximating map
$\ApproximatingMap_{\FamilyIndex}(z)$
by taking the product of each value
$\InputMap_{\FamilyIndex}(\DomainPoint)$
with finite equilateral spaces of cardinality
$\EquilateralSize_{\FamilyIndex}$.
We then prove that the resulting family of images is locally finite and discrete.

\ProofStep{Product approximations.}\label{step:discrete-products}
\paraid{p-021b6255b9ed}
By \Cref{lem:cover-adapted-scale},
choose
$\CoverScale\colon\GHSpace\to(0,1/16]$
satisfying \eqref{eq:scale-lipschitz-2} and the cover condition in that lemma.

\paraid{p-71bd13b39f7a}
For an integer
$\EquilateralSize\geq2$,
let
$\EquilateralCarrier_{\EquilateralSize}=\{0,\ldots,\EquilateralSize-1\}$.
For
$\ProductScale>0$,
let
$\EquilateralMetric_{\EquilateralSize,\ProductScale}$
be the equilateral metric on
$\EquilateralCarrier_{\EquilateralSize}$
with off-diagonal distance
$\ProductScale$.
For each integer
$\EquilateralSize\geq2$,
define
$\ProductApproximation_{\EquilateralSize}\colon\GHSpace\to\GHSpace$
as follows.
For
$\MetricSpace{\BaseCarrier}{\MetricSymbol}\in\GHSpace$,
put
\begin{equation}\label{eq:product-approximation-2}
 \ProductScale=\CoverScale\MetricSpace{\BaseCarrier}{\MetricSymbol}
\end{equation}
and
\begin{equation}\label{eq:product-approximation-1}
 \ProductApproximation_{\EquilateralSize}\MetricSpace{\BaseCarrier}{\MetricSymbol}
 =\MetricSpace{\BaseCarrier\times\EquilateralCarrier_{\EquilateralSize}}
              {\MetricSymbol\vee\EquilateralMetric_{\EquilateralSize,\ProductScale}}.
\end{equation}
The graph of the projection onto
$\BaseCarrier$
is a correspondence of distortion at most
$\ProductScale$.
Hence
we have
\begin{equation}\label{eq:product-error}
 \GHDistance(\ProductApproximation_{\EquilateralSize}\MetricSpace{\BaseCarrier}{\MetricSymbol},
             \MetricSpace{\BaseCarrier}{\MetricSymbol})
 \leq\frac12\CoverScale\MetricSpace{\BaseCarrier}{\MetricSymbol}.
\end{equation}

\paraid{p-7490df62137b}
For an integer
$\EquilateralSize\geq2$
and spaces
$\MetricSpace{\BaseCarrier}{\MetricSymbol},
 \MetricSpace{\ComparisonCarrier}{\ComparisonMetric}\in\GHSpace$,
pair equal labels in the two equilateral factors.
As in the proof of
\cite[Proposition~5.3]{Ishiki2022Branching},
the product correspondence and \eqref{eq:scale-lipschitz-2} imply
\begin{equation}\label{eq:equilateral-product-lipschitz}
 \begin{aligned}
 &\GHDistance(\ProductApproximation_{\EquilateralSize}\MetricSpace{\BaseCarrier}{\MetricSymbol},
              \ProductApproximation_{\EquilateralSize}\MetricSpace{\ComparisonCarrier}{\ComparisonMetric})\\
 &\quad\leq\max\left\{
     \GHDistance(\MetricSpace{\BaseCarrier}{\MetricSymbol},
                 \MetricSpace{\ComparisonCarrier}{\ComparisonMetric}),
     \frac12\left|\CoverScale\MetricSpace{\BaseCarrier}{\MetricSymbol}
                 -\CoverScale\MetricSpace{\ComparisonCarrier}{\ComparisonMetric}\right|
                  \right\}\\
 &\quad\leq\GHDistance(\MetricSpace{\BaseCarrier}{\MetricSymbol},
                      \MetricSpace{\ComparisonCarrier}{\ComparisonMetric}).
 \end{aligned}
\end{equation}
Thus every map
$\ProductApproximation_{\EquilateralSize}$
is
$1$-Lipschitz.

\ProofStep{Choice of the cardinalities and separation of the images.}\label{step:discrete-separation}
\paraid{p-04700bc9ef5a}
We choose the positive integers
$\EquilateralSize_{\FamilyIndex}$
and the maps
$\ApproximatingMap_{\FamilyIndex}$
recursively over
$\FamilyIndex\in\DomainIndexSet$.
Fix
$\FamilyIndex\in\DomainIndexSet$
and assume that the maps with smaller indices have been constructed
and have compact images.
Put
\begin{equation}\label{eq:previous-approximation-images}
 \PreviousApproximationImages_{\FamilyIndex}
 =\bigcup_{\substack{k\in\DomainIndexSet\\k<\FamilyIndex}}
       \ApproximatingMap_k(\ApproximationDomain_k).
\end{equation}
This set is compact,
and
$\PreviousApproximationImages_0=\emptyset$.
Since
$\ApproximationDomain_{\FamilyIndex}$
is nonempty and compact,
the number
\begin{equation}\label{eq:minimum-product-scale}
 \ProductScale_{\FamilyIndex}
 =\min_{\DomainPoint\in\ApproximationDomain_{\FamilyIndex}}
       \CoverScale(\InputMap_{\FamilyIndex}(\DomainPoint))
\end{equation}
is positive.

\paraid{p-5bd0a1e059c4}
Compactness of
$\PreviousApproximationImages_{\FamilyIndex}$
and Corollary~\ref{cor:uniform-separated-sets}
imply that there is an integer
$\PackingBound_{\FamilyIndex}\geq0$
such that every
$\ProductScale_{\FamilyIndex}/2$-separated
subset of every space in
$\PreviousApproximationImages_{\FamilyIndex}$
has cardinality at most
$\PackingBound_{\FamilyIndex}$.
If
$\PreviousApproximationImages_{\FamilyIndex}=\emptyset$,
take
$\PackingBound_{\FamilyIndex}=0$.

\paraid{p-82b288935ffb}
Choose an integer
$\EquilateralSize_{\FamilyIndex}$
larger than all previously chosen cardinalities and satisfying
\begin{equation}\label{eq:equilateral-cardinality-choice}
 \EquilateralSize_{\FamilyIndex}
 >\max\{\FamilyIndex+1,\PackingBound_{\FamilyIndex}\}.
\end{equation}
Define
$\ApproximatingMap_{\FamilyIndex}\colon\ApproximationDomain_{\FamilyIndex}\to\GHSpace$
by
\begin{equation}\label{eq:discrete-maps}
 \ApproximatingMap_{\FamilyIndex}
 =\ProductApproximation_{\EquilateralSize_{\FamilyIndex}}\circ\InputMap_{\FamilyIndex}.
\end{equation}
The map
$\ApproximatingMap_{\FamilyIndex}$
is continuous by \eqref{eq:equilateral-product-lipschitz},
and its image is compact.
For every
$\DomainPoint\in\ApproximationDomain_{\FamilyIndex}$,
equation \eqref{eq:product-error} implies
\begin{equation}\label{eq:approximation-error}
 \GHDistance(\ApproximatingMap_{\FamilyIndex}(\DomainPoint),
             \InputMap_{\FamilyIndex}(\DomainPoint))
 \leq\frac12\CoverScale(\InputMap_{\FamilyIndex}(\DomainPoint))
 <\CoverScale(\InputMap_{\FamilyIndex}(\DomainPoint)).
\end{equation}
Thus
$\ApproximatingMap_{\FamilyIndex}$
and
$\InputMap_{\FamilyIndex}$
are
$\OpenCover$-close.

\paraid{p-dbde286ff7a5}
Every product in
$\ApproximatingMap_{\FamilyIndex}(\ApproximationDomain_{\FamilyIndex})$
contains a copy of its equilateral factor,
obtained by fixing a point in its first factor.
By \eqref{eq:minimum-product-scale},
this copy has
$\EquilateralSize_{\FamilyIndex}$
points separated by at least
$\ProductScale_{\FamilyIndex}$.
If the GH distance from such a product to a space in
$\PreviousApproximationImages_{\FamilyIndex}$
were less than
$\ProductScale_{\FamilyIndex}/4$,
a correspondence of distortion less than
$\ProductScale_{\FamilyIndex}/2$
would send these points to distinct points separated by more than
$\ProductScale_{\FamilyIndex}/2$
in that space.
This would imply
$\EquilateralSize_{\FamilyIndex}\leq\PackingBound_{\FamilyIndex}$,
contrary to \eqref{eq:equilateral-cardinality-choice}.
Consequently,
for every
$\DomainPoint\in\ApproximationDomain_{\FamilyIndex}$
and every
$\MetricSpace{\ComparisonCarrier}{\ComparisonMetric}\in\PreviousApproximationImages_{\FamilyIndex}$,
we have
\begin{equation}\label{eq:recursive-image-separation}
 \GHDistance(\ApproximatingMap_{\FamilyIndex}(\DomainPoint),
             \MetricSpace{\ComparisonCarrier}{\ComparisonMetric})
 \geq\frac{\ProductScale_{\FamilyIndex}}4.
\end{equation}
This completes the recursive construction and proves that the images are pairwise disjoint.

\paraid{p-4a46362ad0ff}
Figure~\ref{fig:discrete-approximation} shows how the scales and cardinalities
control the approximation and the separation of the images.
\begin{figure}[H]
\centering
\begin{tikzpicture}[>=Stealth,font=\small]
\node at (0,1.6) {$\BaseCarrier$};
\draw (0,0) ellipse (1.15 and 0.4);
\fill (-0.55,0.04) circle (1.1pt);
\fill (0.5,-0.12) circle (1.1pt);
\fill (0.05,0.08) circle (1.5pt);
\node[above,inner sep=3pt] at (0.05,0.08) {$\BasePoint$};
\draw[->] (1.5,0) -- node[above]
 {$\times\EquilateralCarrier_3$} (3.1,0);
\node at (4.65,1.6) {$\BaseCarrier\times\EquilateralCarrier_3$};
\fill[black!8] (4.5,-1.15) rectangle (4.9,1.3);
\foreach \level/\index in {0.95/0,0/1,-0.95/2}
{
 \draw (4.65,\level) ellipse (1.15 and 0.4);
 \fill (4.1,{\level+0.04}) circle (1.1pt);
 \fill (5.15,{\level-0.12}) circle (1.1pt);
 \fill (4.7,{\level+0.08}) circle (1.5pt);
 \node[right] at (5.95,\level)
  {$\BaseCarrier\times\{\index\}$};
}
\node[align=center] at (4.7,-1.95)
 {$\{\BasePoint\}\times\EquilateralCarrier_3$\\
  pairwise distance $\ProductScale$};
\end{tikzpicture}
\par\medskip
\begin{tikzpicture}[
 >=Stealth,
 font=\small,
 role/.style={draw,align=center,text width=3.4cm,
              minimum height=2.1cm,inner sep=5pt,anchor=north}
]
\node at (4,0.9)
 {$\MetricSpace{\ClosedDomain}{\SubspaceMetric}
   =\ProductApproximation_{\EquilateralSize_{\FamilyIndex}}
      \MetricSpace{\BaseCarrier}{\MetricSymbol},
   \qquad
   \ProductScale=\CoverScale\MetricSpace{\BaseCarrier}{\MetricSymbol}$};
\node (formula) at (4,0)
 {$\MetricSpace{\ClosedDomain}{\SubspaceMetric}
   =\MetricSpace{\BaseCarrier\times\EquilateralCarrier_{\EquilateralSize_{\FamilyIndex}}}
                {\MetricSymbol\vee\EquilateralMetric_{\EquilateralSize_{\FamilyIndex},\ProductScale}}$};
\node[role] (approximation) at (0,-1.7)
 {small scale $\ProductScale$\\
  $\begin{gathered}
    \GHDistance(\MetricSpace{\ClosedDomain}{\SubspaceMetric},
                 \MetricSpace{\BaseCarrier}{\MetricSymbol})\\
    \leq\ProductScale/2
   \end{gathered}$\\
  approximation};
\node[role] (separation) at (4,-1.7)
 {$\EquilateralSize_{\FamilyIndex}>\PackingBound_{\FamilyIndex}$\\
  recursive choice\\
  disjoint images};
\node[role] (packing) at (8,-1.7)
 {$\EquilateralSize_{\FamilyIndex}\to\infty$\\
  positive scale bound\\
  no accumulation};
\draw[->,shorten <=2pt]
 ([xshift=-1.8cm]formula.south) .. controls +(0,-0.55) and +(0,0.55) .. (approximation.north);
\draw[->,shorten <=2pt]
 (formula.south) -- (separation.north);
\draw[->,shorten <=2pt]
 ([xshift=1.8cm]formula.south) .. controls +(0,-0.55) and +(0,0.55) .. (packing.north);
\end{tikzpicture}
\caption{The upper diagram shows three copies of
$\BaseCarrier$
in a product with an equilateral space.
The marked fiber has
$\EquilateralSize$
points at pairwise distance
$\ProductScale$
for a general factor
$\EquilateralCarrier_\EquilateralSize$.
The scales and cardinalities of the equilateral factors control the
 discrete approximation.
The bound
$\PackingBound_{\FamilyIndex}$
applies to subsets separated by
$\ProductScale_{\FamilyIndex}/2$
in the previously constructed images,
where
$\ProductScale_{\FamilyIndex}$
is the minimum input scale in \eqref{eq:minimum-product-scale}.
The recursive choice \eqref{eq:equilateral-cardinality-choice}
separates the images.
Along a convergent sequence of approximating spaces,
the scales have a positive lower bound,
so the divergence of the cardinalities contradicts
Corollary~\ref{cor:uniform-separated-sets}.}
\label{fig:discrete-approximation}
\end{figure}


\ProofStep{Local finiteness.}\label{step:discrete-local-finiteness}
\paraid{p-86b6e1b5bb9e}
For the sake of contradiction,
suppose that the family
$\{\ApproximatingMap_{\FamilyIndex}(\ApproximationDomain_{\FamilyIndex})\}_{\FamilyIndex\in\DomainIndexSet}$
is not locally finite at some
$\MetricSpace{\LimitCarrier}{\LimitMetric}\in\GHSpace$.
Then every neighborhood of this point meets infinitely many images.
For each
$\SequenceIndex\in\NonnegativeIntegers$,
choose an index
$\FamilyIndex_{\SequenceIndex}$
larger than all previously chosen indices and a point
$\DomainPoint_{\SequenceIndex}\in\ApproximationDomain_{\FamilyIndex_{\SequenceIndex}}$
whose image lies in the GH ball of radius
$2^{-\SequenceIndex}$
about
$\MetricSpace{\LimitCarrier}{\LimitMetric}$.
For each
$\SequenceIndex\in\NonnegativeIntegers$,
put
$\InputMap_{\FamilyIndex_{\SequenceIndex}}(\DomainPoint_{\SequenceIndex})
=\MetricSpace{\BaseCarrier_{\SequenceIndex}}{\MetricSymbol_{\SequenceIndex}}$,
and
$\ApproximatingMap_{\FamilyIndex_{\SequenceIndex}}(\DomainPoint_{\SequenceIndex})
=\MetricSpace{\ComparisonCarrier_{\SequenceIndex}}{\ComparisonMetric_{\SequenceIndex}}$,
given by \eqref{eq:product-approximation-1} and \eqref{eq:discrete-maps}.
Then
\begin{equation}\label{eq:equilateral-image-convergence}
 \MetricSpace{\ComparisonCarrier_{\SequenceIndex}}{\ComparisonMetric_{\SequenceIndex}}
 \to\MetricSpace{\LimitCarrier}{\LimitMetric}.
\end{equation}
For each
$\SequenceIndex\in\NonnegativeIntegers$,
put
$\CoverScale_{\SequenceIndex}
 =\CoverScale\MetricSpace{\BaseCarrier_{\SequenceIndex}}{\MetricSymbol_{\SequenceIndex}}$.
For every
$\SequenceIndex\in\NonnegativeIntegers$,
equations \eqref{eq:scale-lipschitz-2} and \eqref{eq:approximation-error} imply
\begin{equation}\label{eq:scale-lower-bound}
 \begin{aligned}
 \CoverScale\MetricSpace{\LimitCarrier}{\LimitMetric}
 &\leq\CoverScale_{\SequenceIndex}
      +\frac1{16}\GHDistance(
         \MetricSpace{\BaseCarrier_{\SequenceIndex}}{\MetricSymbol_{\SequenceIndex}},
         \MetricSpace{\LimitCarrier}{\LimitMetric})\\
 &\leq\CoverScale_{\SequenceIndex}
      +\frac1{16}\Bigl(\GHDistance(
         \MetricSpace{\BaseCarrier_{\SequenceIndex}}{\MetricSymbol_{\SequenceIndex}},
         \MetricSpace{\ComparisonCarrier_{\SequenceIndex}}{\ComparisonMetric_{\SequenceIndex}})
         +\GHDistance(
         \MetricSpace{\ComparisonCarrier_{\SequenceIndex}}{\ComparisonMetric_{\SequenceIndex}},
         \MetricSpace{\LimitCarrier}{\LimitMetric})\Bigr)\\
 &\leq\CoverScale_{\SequenceIndex}
      +\frac1{16}\Bigl(\frac12\CoverScale_{\SequenceIndex}
         +\GHDistance(
         \MetricSpace{\ComparisonCarrier_{\SequenceIndex}}{\ComparisonMetric_{\SequenceIndex}},
         \MetricSpace{\LimitCarrier}{\LimitMetric})\Bigr)\\
 &\leq\frac{33}{32}\CoverScale_{\SequenceIndex}
      +\frac1{16}\GHDistance(
         \MetricSpace{\ComparisonCarrier_{\SequenceIndex}}{\ComparisonMetric_{\SequenceIndex}},
         \MetricSpace{\LimitCarrier}{\LimitMetric}).
 \end{aligned}
\end{equation}
Put
$\SeparationScale=\frac12\CoverScale\MetricSpace{\LimitCarrier}{\LimitMetric}>0$.
By \eqref{eq:equilateral-image-convergence} and \eqref{eq:scale-lower-bound},
for all sufficiently large
$\SequenceIndex$
we have
\begin{equation}\label{eq:equilateral-positive-scale}
 \CoverScale_{\SequenceIndex}\geq\SeparationScale.
\end{equation}

\paraid{p-866cf4c6571b}
For every
$\SequenceIndex\in\NonnegativeIntegers$,
choose
$\BasePoint_{\SequenceIndex}\in\BaseCarrier_{\SequenceIndex}$
and define
\[
 \SeparatedSet_{\SequenceIndex}
 =\{(\BasePoint_{\SequenceIndex},\EquilateralPoint)\mid
     \EquilateralPoint\in\EquilateralCarrier_{\EquilateralSize_{\FamilyIndex_{\SequenceIndex}}}\}
 \subset\ComparisonCarrier_{\SequenceIndex}.
\]
For all sufficiently large
$\SequenceIndex$
and all distinct
$\EquilateralPoint_0,\EquilateralPoint_1
 \in\EquilateralCarrier_{\EquilateralSize_{\FamilyIndex_{\SequenceIndex}}}$,
equation \eqref{eq:equilateral-positive-scale} implies
\[
 \ComparisonMetric_{\SequenceIndex}
 ((\BasePoint_{\SequenceIndex},\EquilateralPoint_0),
  (\BasePoint_{\SequenceIndex},\EquilateralPoint_1))
 =\max\{0,\CoverScale_{\SequenceIndex}\}
 =\CoverScale_{\SequenceIndex}\geq\SeparationScale.
\]
Thus
$\SeparatedSet_{\SequenceIndex}$
is
$\SeparationScale$-separated
for all sufficiently large
$\SequenceIndex$.
By \eqref{eq:equilateral-image-convergence}
and Corollary~\ref{cor:uniform-separated-sets},
the cardinalities
$\Card(\SeparatedSet_{\SequenceIndex})$
are uniformly bounded.
This contradicts \eqref{eq:equilateral-cardinality-choice},
since
\[
 \Card(\SeparatedSet_{\SequenceIndex})
 =\EquilateralSize_{\FamilyIndex_{\SequenceIndex}}
 >\FamilyIndex_{\SequenceIndex}+1\longrightarrow\infty.
\]

\paraid{p-eb900edd6e7c}
Therefore the family
$\{\ApproximatingMap_{\FamilyIndex}(\ApproximationDomain_{\FamilyIndex})\}_{\FamilyIndex\in\DomainIndexSet}$
is locally finite.

\ProofStep{Discreteness.}\label{step:discrete-discreteness}
\paraid{p-0f7257809982}
A locally finite family of pairwise disjoint closed sets is discrete.
By Steps~\ref{step:discrete-separation} and~\ref{step:discrete-local-finiteness},
this fact applies to our compact images.
Local finiteness is automatic when
$\DomainIndexSet$
is finite.
\end{proof}
